\section{Active radial interfaces near an initial mark}
\label{sec:area_law_initial_active_rays}

Fix an arbitrary origin $o\in\mathbb R^2$, a finite endpoint set
$Z\subset\mathbb Z^2$, an integer width $C\ge2$ and an initial layer
$k_0\ge50{,}000{,}000$. Choose arbitrary valid pitch residues at every
layer. Let $\mathcal I$ be the actual initial identifiers, let $X_i^0$
be their closed birth regions, and put $U_i=\operatorname{int}(X_i^0)$.
Let $k\ge k_0$, let $z\in F_{k,\ell_k}$ be an actual fine-cell index,
and let $v$ be one of the nine marks of its defining cell.
Put $\ell=\ell_k-5$ and $r=2^\ell/2=t_k/64$.

Join $v$ to the eight elementary bases of the all-midpoint fan in
$\overline B_\infty(v,r)$. Let $\mathcal J$ index those bases, write
$P_s$ for the corresponding closed triangle, and write $b_s^-$ and
$b_s^+$ for its initial and final base endpoints in the perimeter
orientation. Let $\sigma:\mathcal J\to\mathcal I$ be the unique actual
assignment of Theorem~\ref{thm:area_law_actual_sector_assignment}.
Thus, for every $i\in\mathcal I$,
\begin{align}
    X_i^0\cap\overline B_\infty(v,r)
     & =\bigcup_{\substack{s\in\mathcal J \\\sigma(s)=i}}P_s.
    \label{eq:active_radial_decomposition}
\end{align}
An oriented neighboring pair $(s,t)$ is active when
$b_s^+=b_t^-$ and $\sigma(s)\ne\sigma(t)$.
Let $\mathcal A$ be the set of such pairs. This describes the radial
interfaces in the local construction of
\cite[Section~11]{OpenAI2026AreaLaw}.

% Source: 10-geometry.tex, geometry:initial-stars, lines 333--370,
% especially 352--370, and prop:two-families, lines 308--323;
% revision adc7f1241b42e322a6451854ab7e4b4c146bf78a.
% These are the open-radius and closed-half-radius local assertions.
% No full isolated-star theorem or active-ray enumeration is asserted.

\begin{theorem}[Actual frontiers and active radial segments]
    \label{thm:area_law_initial_active_radial}
    \lean{TNLean.PEPS.AreaLaw.Geometry.initialRegion_frontier_near_mark_iff_active_radial}
    \leanok
    \uses{thm:area_law_actual_sector_assignment,
        def:area_law_initial_birth_regions,def:area_law_cell_fan}
    For every $x\in B_\infty(v,r)$,
    \begin{align}
        x\in\bigcup_{i\in\mathcal I}\partial X_i^0
         & \quad\Longleftrightarrow\quad
        \exists(s,t)\in\mathcal A,\quad x\in[v,b_s^+].
    \end{align}
    In the picture below each triangle carries its identifier $\sigma(s)$;
    the solid radial segments are the active ones, and they are exactly the
    frontier points of the birth regions inside the open square.
    % Source: 10-geometry.tex, geometry:initial-stars, lines 333--370, and the
    % left panel of geometry:repair-picture, lines 418--466. Ink-to-index: the
    % red dot is the mark v; the dotted square is the boundary of the closed
    % ball of radius r, with its eight base endpoints (corners and side
    % midpoints) dotted. Each of the eight triangles P_s is labelled by its
    % identifier sigma(s); the sample assignment is i_1,i_1 (upper),
    % i_2,i_2,i_2 (right and lower right), i_3, i_4,i_4. A solid segment
    % [v, b_s^+] separates neighbours with different identifiers (an active
    % pair); a dashed segment separates neighbours with the same identifier and
    % is not a frontier. By the sector-color theorem the colors chi alternate
    % across solid segments only. These lines are geometry, not tensor indices:
    % no contractions, empty boundary signature. Not to scale.
    \begin{center}
        \begin{tenkz}[size=s, metrics=compact, rows={wire,wire,wire,wire,wire,wire,wire,wire,wire}, cols=9, bonds=none]
            \tndeclare{species}{centre}{hue=source:red}
            \tn[at=(5,5), name=starc, species=centre, ports={0:virtual, 45:virtual, 90:virtual, 135:virtual, 180:virtual, 225:virtual, 270:virtual, 315:virtual}]{}
            \tn[at=(1,1), name=starnw, ports={0:virtual, 270:virtual, 315:virtual}]{}
            \tn[at=(1,5), name=starn, ports={0:virtual, 180:virtual, 270:virtual}]{}
            \tn[at=(1,9), name=starne, ports={180:virtual, 225:virtual, 270:virtual}]{}
            \tn[at=(5,9), name=stare, ports={90:virtual, 180:virtual, 270:virtual}]{}
            \tn[at=(9,9), name=starse, ports={90:virtual, 135:virtual, 180:virtual}]{}
            \tn[at=(9,5), name=stars, ports={0:virtual, 90:virtual, 180:virtual}]{}
            \tn[at=(9,1), name=starsw, ports={0:virtual, 45:virtual, 90:virtual}]{}
            \tn[at=(5,1), name=starw, ports={0:virtual, 90:virtual, 270:virtual}]{}
            \tnwire[stroke=dotted, name=starsidenw]{starnw.0}{starn.180}
            \tnwire[stroke=dotted, name=starsiden]{starn.0}{starne.180}
            \tnwire[stroke=dotted, name=starsidene]{starne.270}{stare.90}
            \tnwire[stroke=dotted, name=starsidee]{stare.270}{starse.90}
            \tnwire[stroke=dotted, name=starsidese]{starse.180}{stars.0}
            \tnwire[stroke=dotted, name=starsides]{stars.180}{starsw.0}
            \tnwire[stroke=dotted, name=starsidesw]{starsw.90}{starw.270}
            \tnwire[stroke=dotted, name=starsidew]{starw.90}{starnw.270}
            \tnwire[name=starraynw]{starc.135}{starnw.315}
            \tnwire[stroke=dashed, name=starrayn]{starc.90}{starn.270}
            \tnwire[name=starrayne]{starc.45}{starne.225}
            \tnwire[stroke=dashed, name=starraye]{starc.0}{stare.180}
            \tnwire[stroke=dashed, name=starrayse]{starc.315}{starse.135}
            \tnwire[name=starrays]{starc.270}{stars.90}
            \tnwire[name=starraysw]{starc.225}{starsw.45}
            \tnwire[stroke=dashed, name=starrayw]{starc.180}{starw.0}
            \tnmark[form=label, label pos=22.5]{(5,5)}{$v$}
            \tnmark[form=label]{(3,4)}{$i_1$}
            \tnmark[form=label]{(3,6)}{$i_1$}
            \tnmark[form=label]{(4,8)}{$i_2$}
            \tnmark[form=label]{(7,8)}{$i_2$}
            \tnmark[form=label]{(8,6)}{$i_2$}
            \tnmark[form=label]{(8,4)}{$i_3$}
            \tnmark[form=label]{(7,2)}{$i_4$}
            \tnmark[form=label]{(4,2)}{$i_4$}
            \tnmark[form=label, label pos=se]{(9,9)}{$\overline B_\infty(v,r)$}
        \end{tenkz}
    \end{center}
\end{theorem}
\begin{proof}\leanok
    \uses{thm:area_law_actual_sector_assignment,
        thm:area_law_finite_cover_frontier,
        thm:area_law_cell_fan_regularity,
        thm:area_law_cell_fan_cover,
        thm:area_law_centered_dyadic_cell_ball,
        thm:area_law_initial_birth_regularity,
        thm:area_law_initial_open_disjoint,
        thm:area_law_initial_sector_no_active,
        thm:area_law_fan_triangle_contact,
        thm:area_law_cell_fan_intersections}
    Put $K=\sigma(\mathcal J)$ and, for each $i\in K$, put
    $E_i=\bigcup_{\sigma(s)=i}P_s$.
    This is a finite closed family covering
    $\overline B_\infty(v,r)$.
    By~(\ref{eq:active_radial_decomposition}), $X_i^0$ and $E_i$
    agree on a neighborhood of every point of $B_\infty(v,r)$.
    Their interior membership therefore agrees there. An identifier
    outside $K$ has no birth-region point in this neighborhood.

    Suppose that $x\in\partial X_i^0$ and $x\ne v$.
    The finite-cover frontier theorem gives a different identifier
    $j\in K$ with $x\in E_i\cap E_j$; its outer-frontier alternative
    is excluded by $x\in B_\infty(v,r)$.
    Choose $s,t\in\mathcal J$ with $\sigma(s)=i$, $\sigma(t)=j$
    and $x\in P_s\cap P_t$.
    Both triangles also contain $v$, so their intersection is
    nontrivial. The actual triangle-contact theorem gives, after
    interchanging $s$ and $t$ if necessary,
    $b_s^+=b_t^-$ and $P_s\cap P_t=[v,b_s^+]$.
    Since $i\ne j$, this pair is active.

    If $x=v$ and no pair is active, every adjacent pair has the same
    assigned identifier. The no-active theorem then places the smaller
    closed square in one actual open region $U_j$.
    A point of $U_j$ is on no birth frontier: it is interior to $X_j^0$,
    and disjointness of the open regions, together with
    $X_i^0=\overline{U_i}$, excludes membership in $X_i^0$ for $i\ne j$.
    Thus a frontier point at $v$ forces an active pair, whose radial
    segment contains $v$.

    Conversely, let $(s,t)\in\mathcal A$ and $x\in[v,b_s^+]$.
    The exact fan-intersection formula includes this segment in
    $P_s\cap P_t$: the common base endpoint belongs to both bases,
    and its radial join with $v$ belongs to both triangles.
    By~(\ref{eq:active_radial_decomposition}),
    $x\in X_{\sigma(s)}^0\cap X_{\sigma(t)}^0$.
    Since the two identifiers are different,
    $U_{\sigma(s)}$ is disjoint from
    $X_{\sigma(t)}^0=\overline{U_{\sigma(t)}}$.
    Hence $x\notin U_{\sigma(s)}$ and
    $x\in\partial X_{\sigma(s)}^0$, as required.
\end{proof}

\begin{theorem}[Active radial description on a smaller closed square]
    \label{thm:area_law_initial_active_radial_closed}
    \lean{TNLean.PEPS.AreaLaw.Geometry.initialRegion_frontier_small_closedBall_iff_active_radial}
    \leanok
    \uses{thm:area_law_actual_sector_assignment,
        def:area_law_initial_birth_regions,def:area_law_cell_fan}
    For every $x\in\overline B_\infty(v,r/2)$,
    \begin{align}
        x\in\bigcup_{i\in\mathcal I}\partial X_i^0
         & \quad\Longleftrightarrow\quad
        \exists(s,t)\in\mathcal A,\quad x\in[v,b_s^+].
    \end{align}
    The closed half-side is $r/2=t_k/128$. If the reference cell
    realizes the smallest incident side $S_v=t_k$, it is $S_v/128$.
\end{theorem}
\begin{proof}\leanok
    \uses{thm:area_law_initial_active_radial}
    The fan radius is positive. Consequently,
    $\overline B_\infty(v,r/2)\subseteq B_\infty(v,r)$, and the
    open-square equivalence applies to every point of this closed square.
\end{proof}
